% Introduction (issue #7).

\section{Introduction}
\label{sec:introduction}

An aircraft engine that fails in service grounds the aircraft and can endanger it,
while an engine that is overhauled too early wastes useful life. Predictive
maintenance resolves this trade-off by estimating, from the engine's own sensor
history, its remaining useful life (RUL): the number of operating cycles left
before failure~\cite{Jardine2006review,Lei2018machinery}. Because run-to-failure
records of real engines are rare, most data-driven work uses NASA's C-MAPSS
simulation data~\cite{Saxena2008damage,Saxena2008turbofan}, in which every
training engine is observed until it fails and every test engine stops at an
unknown point before failure.

This project asks whether hand-engineered degradation features are enough for RUL
prediction, or whether models that learn directly from raw sensor sequences do
better. Fixed rolling summaries of recent history may miss long degradation
trajectories and interactions whose timing differs from engine to engine, which
sequence models are designed to capture~\cite{Zheng2017long}. Answering the
question fairly needs a strong feature-based baseline and an evaluation protocol
that both model families share. Published C-MAPSS results are hard to compare
because they differ in subsets, targets and protocols~\cite{Ramasso2014performance}.

Phase~1, reported here, builds that baseline on FD001, the subset with one
operating condition and one fault mode. Its contributions are:
\begin{itemize}
\item a leakage-controlled protocol, fixed for Phase~2: engine-grouped
  cross-validation, preprocessing refitted inside every training fold, an
  uncapped RUL target, and scoring of each official test engine once at its last
  cycle with RMSE, MAE and the asymmetric NASA score, overall and by degradation
  stage;
\item a causal feature set of 155 predictors per cycle (regime-normalized sensor
  readings with rolling means, slopes and changes from each engine's first
  cycle), compared with 25 raw inputs;
\item Random Forest, XGBoost and Elastic Net models tuned under this protocol and
  scored on the 100 official FD001 test engines;
\item an error analysis showing that the engineered features roughly halve the
  tree models' error near failure but overestimate young engines, that the NASA
  score is dominated by single engines, and that capping the target reverses the
  raw-versus-engineered comparison.
\end{itemize}

Section~\ref{sec:literature_review} reviews related work and
Section~\ref{sec:dataset_eda} the data. Section~\ref{sec:methodology} defines the
problem and the protocol, Sections~\ref{sec:feature_engineering}
and~\ref{sec:models} describe the features and models, and
Section~\ref{sec:results} reports the results. Section~\ref{sec:conclusion}
concludes and outlines Phase~2.
